\documentclass[11pt,a4paper]{article}
\usepackage[margin=2.1cm]{geometry}\usepackage{fontspec}\usepackage[italian]{babel}\usepackage{xcolor,longtable,booktabs,hyperref,fancyhdr}
\setmainfont{NimbusSans-Regular.otf}[Path=sorgenti/font/,BoldFont=NimbusSans-Bold.otf,ItalicFont=NimbusSans-Italic.otf]
\setmonofont{lmmono10-regular.otf}[Path=sorgenti/font/,Scale=.88]
\definecolor{teal}{HTML}{087C74}\hypersetup{colorlinks=true,linkcolor=teal,urlcolor=teal}\setlength{\headheight}{14pt}\setlength{\parindent}{0pt}\setlength{\parskip}{7pt}\pagestyle{fancy}\fancyhf{}\fancyfoot[R]{\thepage}\fancyhead[L]{Geologia · Statistica e coding R · R6}\begin{document}


\section*{M05 · Laboratori e soluzioni}

Le venti attività corrispondono ai file M05\_N01.R–M05\_N20.R. Ogni script ricostruisce i calcoli condivisi e stampa i risultati pertinenti; viene eseguito dalla radice del progetto. Le pagine indicate si riferiscono al fascicolo M05 della revisione approvata, mantenuto a 105 slide.

\newpage

\subsection*{N01 · Parametro, campione e statistica}

\textbf{Pagine:} 4–6

\textbf{Obiettivo.} Distinguere il parametro dal valore osservato dello stimatore.

\textbf{Consegna.} Eseguire il campionamento esponenziale, riportare parametro, numerosità e media osservata. Ripetere con un altro seme, spiegando quale oggetto cambi e quale rimanga fissato.

\textbf{Soluzione e interpretazione.} Il parametro è la media teorica 1 della distribuzione esponenziale. La numerosità è 40 e la stima varia con il campione. Il cambiamento del seme seleziona un’altra realizzazione del procedimento; non modifica il parametro del modello.

\textbf{Variazione.} Confrontare media e mediana dello stesso campione, indicando i rispettivi parametri.

\texttt{laboratori/M05/M05\_N01.R}

\subsection*{N02 · Media cumulativa e convergenza}

\textbf{Pagine:} 30

\textbf{Obiettivo.} Interpretare la convergenza senza attribuirle monotonicità.

\textbf{Consegna.} Rappresentare i quattro percorsi della media cumulativa e individuare tratti nei quali una nuova osservazione allontani la media da 1. Esplicitare la disuguaglianza di Chebyshev utilizzata nelle slide.

\textbf{Soluzione e interpretazione.} La convergenza in probabilità riguarda la probabilità di uno scostamento superiore a una tolleranza, quando n aumenta. Un singolo percorso può oscillare e allontanarsi temporaneamente. Il limite 1/(n epsilon\textsuperscript{2}) deriva dalla varianza unitaria e dall’indipendenza.

\textbf{Variazione.} Modificare la distribuzione esponenziale con rate=2 e aggiornare il parametro di riferimento.

\texttt{laboratori/M05/M05\_N02.R}

\newpage

\subsection*{N03 · Distribuzione delle medie}

\textbf{Pagine:} 22–26

\textbf{Obiettivo.} Identificare la media come unità rappresentata nell’istogramma.

\textbf{Consegna.} Confrontare le 6000 medie per n=5 e n=40. Calcolarne media e deviazione standard e confrontarle con 1 e 1/sqrt(n).

\textbf{Soluzione e interpretazione.} Le distribuzioni delle medie sono centrate vicino a 1. Le deviazioni standard simulate sono 0,440 e 0,158, vicine ai valori teorici 0,447 e 0,158. Le differenze residue riflettono una simulazione finita.

\textbf{Variazione.} Aggiungere n=15 mantenendo lo stesso numero di simulazioni.

\texttt{laboratori/M05/M05\_N03.R}

\subsection*{N04 · Standardizzazione e dipendenza}

\textbf{Pagine:} 27–33

\textbf{Obiettivo.} Collegare il limite centrale alla scala della statistica e la dipendenza alla sua varianza.

\textbf{Consegna.} Calcolare centro e dispersione delle medie standardizzate. Ricostruire la varianza di una media con sei blocchi indipendenti di quattro unità e ICC=0,6, usando sia la formula sia la matrice Sigma.

\textbf{Soluzione e interpretazione.} Le medie standardizzate hanno centro vicino a zero e dispersione vicina a uno. Nel modello a blocchi, la varianza è 0,1166667, DEFF è 2,8 e n effettivo è 8,5714. I due calcoli della varianza coincidono; le ipotesi di equicorrelazione e indipendenza fra blocchi sono essenziali.

\textbf{Variazione.} Sostituire ICC=0,3 e confrontare la variazione di n effettivo.

\texttt{laboratori/M05/M05\_N04.R}

\newpage

\subsection*{N05 · Dispersione dei dati ed errore standard}

\textbf{Pagine:} 34–38, 43

\textbf{Obiettivo.} Calcolare e interpretare SD, SE e barre di incertezza.

\textbf{Consegna.} Per Rounded calcolare numerosità, media, SD e SE. Spiegare quale quantità si dimezzi se la numerosità indipendente quadruplica a varianza invariata.

\textbf{Soluzione e interpretazione.} La media è 6,064 µm, SD è 1,000 µm e SE è 0,112 µm. Il raddoppio della radice della numerosità dimezza SE, mentre la variabilità individuale resta descritta da sigma. Le barre ± SD, ± SE e IC hanno significati differenti.

\textbf{Variazione.} Calcolare l’errore standard teorico di una proporzione per n=40 e p=0,1; discutere il caso di zero successi osservati.

\texttt{laboratori/M05/M05\_N05.R}

\subsection*{N06 · Intervallo di Student}

\textbf{Pagine:} 44–48, 50–52

\textbf{Obiettivo.} Derivare e verificare i limiti dell’intervallo per una media.

\textbf{Consegna.} Calcolare il quantile critico, la semiampiezza e i limiti per Rounded. Confrontare il risultato con t.test(x) e distinguere l’intervallo per la media da quello predittivo.

\textbf{Soluzione e interpretazione.} Il quantile t è 1,990 e l’IC per la media è [5,841; 6,286] µm. L’intervallo predittivo sotto normalità IID è [4,060; 8,067] µm, poiché include la variabilità della nuova osservazione. Il test numerico confronta i limiti calcolati esplicitamente con quelli di R.

\textbf{Variazione.} Passare dal livello 95\% al 99\% e spiegare l’aumento della semiampiezza.

\texttt{laboratori/M05/M05\_N06.R}

\newpage

\subsection*{N07 · Copertura della procedura}

\textbf{Pagine:} 49, 97–98

\textbf{Obiettivo.} Interpretare una probabilità di copertura e il suo errore Monte Carlo.

\textbf{Consegna.} Calcolare la frazione di intervalli contenenti la media 5 nei 2000 campioni. Confrontarla con 0,95 e con l’errore Monte Carlo teorico.

\textbf{Soluzione e interpretazione.} In questa esecuzione 1900 intervalli su 2000 contengono il parametro. La copertura empirica è 0,95 e il SE Monte Carlo teorico è circa 0,00487. L’uguaglianza fra valore empirico e nominale è un esito di questa simulazione, non una necessità.

\textbf{Variazione.} Generare dati esponenziali e confrontare la copertura t con quella dell’esperimento normale, senza attribuirle validità esatta.

\texttt{laboratori/M05/M05\_N07.R}

\subsection*{N08 · Precisione, selezione e popolazione finita}

\textbf{Pagine:} 12, 15–21, 39

\textbf{Obiettivo.} Separare precisione e distorsione e verificare una varianza di disegno.

\textbf{Consegna.} Confrontare i due scenari di selezione per n=10,40,160. Enumerare inoltre tutti i campioni di due unità dalla popolazione 2,4,6,8,10 e confrontare la varianza delle medie con la correzione finita.

\textbf{Soluzione e interpretazione.} La selezione centrata sul valore 1 resta distorta rispetto al parametro territoriale 0, anche quando la distribuzione si restringe. Nell’esempio finito esistono dieci campioni equiprobabili; la varianza esatta delle loro medie è 3, uguale alla formula con fattore 1−n/N.

\textbf{Variazione.} Portare la numerosità del campione finito a tre e ripetere l’enumerazione.

\texttt{laboratori/M05/M05\_N08.R}

\newpage

\subsection*{N09 · Pesi, precisione richiesta e conclusione}

\textbf{Pagine:} 9–10, 40–41, 92–94

\textbf{Obiettivo.} Motivare ponderazione e numerosità in relazione all’estimando.

\textbf{Consegna.} Confrontare la media equiponderata delle stazioni con lo scenario che assegna peso quattro ai provini A1. Calcolare n per sigma pilota 0,8 µm e semiampiezza desiderata 0,2 µm al 95\%.

\textbf{Soluzione e interpretazione.} La media equiponderata è 2,599 g/cm\textsuperscript{3} e quella ponderata è 2,628 g/cm\textsuperscript{3}. La numerosità orientativa è 62 unità indipendenti. La formula di pianificazione non incorpora automaticamente dipendenza, perdite o distorsione di selezione.

\textbf{Variazione.} Risolvere l’esercizio delle due stazioni con numerosità 2 e 8 e medie 2,4 e 2,8.

\texttt{laboratori/M05/M05\_N09.R}

\subsection*{N10 · Distribuzione empirica e informazione}

\textbf{Pagine:} 55–56, 85

\textbf{Obiettivo.} Descrivere l’oggetto condizionato costruito dal bootstrap.

\textbf{Consegna.} Identificare supporto e probabilità della distribuzione empirica Rounded. Calcolare il numero atteso di individui distinti in una replica di ottanta estrazioni.

\textbf{Soluzione e interpretazione.} Ogni indice ha probabilità 1/80 a ogni estrazione. La probabilità che un individuo compaia almeno una volta è 1−(79/80)\textasciicircum{}80, per cui il numero atteso di distinti è circa 50,75. Il numero delle estrazioni resta 80 e il numero di nuove osservazioni reali è zero.

\textbf{Variazione.} Confrontare la proporzione attesa di distinti per n=4,20,80 e 1000.

\texttt{laboratori/M05/M05\_N10.R}

\newpage

\subsection*{N11 · Una replica osservata negli indici}

\textbf{Pagine:} 56–58

\textbf{Obiettivo.} Ricostruire una replica con indici e molteplicità equivalenti.

\textbf{Consegna.} Con il seme 311 estrarre ottanta indici. Contare gli individui distinti e verificare che la media dei valori selezionati coincida con la media ponderata dalle molteplicità.

\textbf{Soluzione e interpretazione.} La replica contiene 55 individui distinti e ha media 6,229 µm. La somma delle molteplicità è 80; la formula sum(frequenze*x)/80 coincide con mean(replica). L’equivalenza permette di controllare che indice, unità e valore rimangano collegati.

\textbf{Variazione.} Estrarre una nuova replica e individuare i cinque indici con maggiore molteplicità.

\texttt{laboratori/M05/M05\_N11.R}

\subsection*{N12 · Bootstrap della media di un gruppo}

\textbf{Pagine:} 61–66

\textbf{Obiettivo.} Confrontare simulazione e momenti condizionati esatti.

\textbf{Consegna.} Calcolare media delle repliche, SE bootstrap e quantili percentili per Rounded. Confrontare SE simulato con sqrt((n−1)*var(x)/n\textasciicircum{}2).

\textbf{Soluzione e interpretazione.} Il SE simulato è 0,112 µm e quello condizionato esatto è circa 0,11115 µm. I limiti percentili sono [5,843; 6,286] µm. Per la media, il valore atteso condizionato è esattamente la media osservata; un piccolo scarto numerico fra le due medie deriva da B finito.

\textbf{Variazione.} Usare B=500 e confrontare la precisione numerica con B=10000.

\texttt{laboratori/M05/M05\_N12.R}

\newpage

\subsection*{N13 · Pooled e quote casuali}

\textbf{Pagine:} 69, 72

\textbf{Obiettivo.} Riconoscere la variabilità della composizione nel ricampionamento pooled.

\textbf{Consegna.} Calcolare la media e la deviazione standard del numero di Spiky nelle repliche. Confrontarle con 80 e sqrt(40), quindi riportare l’IC della media pooled.

\textbf{Soluzione e interpretazione.} Nel pooled il conteggio Spiky segue condizionatamente una binomiale con 160 estrazioni e probabilità 0,5. La media pooled osservata è 6,355 µm; SE è 0,079 µm e l’IC percentile è [6,201; 6,506] µm. Le quote fluttuano fra repliche.

\textbf{Variazione.} Separare il contributo entro morfotipo da quello fra morfotipi usando la varianza empirica con denominatore n.

\texttt{laboratori/M05/M05\_N13.R}

\subsection*{N14 · Bootstrap stratificato}

\textbf{Pagine:} 70–73

\textbf{Obiettivo.} Conservare le quote del generatore e i pesi dell’estimando.

\textbf{Consegna.} Eseguire il bootstrap entro i due gruppi e confrontare SE e intervallo con quelli pooled. Esplicitare perché la media osservata coincida.

\textbf{Soluzione e interpretazione.} Le quote sono fissate a 80/80 e i pesi a 1/2. Il SE stratificato è 0,077 µm e l’IC è [6,201; 6,504] µm. La media coincide con quella pooled perché il campione osservato è bilanciato; le repliche seguono però disegni diversi.

\textbf{Variazione.} Definire un estimando con pesi 0,3 e 0,7, mantenendo esplicita la differenza fra pesi e quote campionarie.

\texttt{laboratori/M05/M05\_N14.R}

\newpage

\subsection*{N15 · Contrasti indipendenti e coppie}

\textbf{Pagine:} 51–52, 74–75

\textbf{Obiettivo.} Mantenere l’orientamento del contrasto e l’appaiamento.

\textbf{Consegna.} Riportare il bootstrap del contrasto Spiky meno Rounded. Per Rounded, calcolare inoltre la differenza totale meno interno e ricampionare le differenze individuali.

\textbf{Soluzione e interpretazione.} Il contrasto fra morfotipi è 0,582 µm con IC [0,286; 0,881] µm. La differenza appaiata Rounded è 0,857 µm, con SE bootstrap 0,027 µm. L’indipendenza fra gruppi e la dipendenza entro coppia richiedono algoritmi diversi.

\textbf{Variazione.} Calcolare la differenza fra le varianze ottenute conservando o eliminando il termine di covarianza entro coppia.

\texttt{laboratori/M05/M05\_N15.R}

\subsection*{N16 · Provini e medie di stazione}

\textbf{Pagine:} 77–79, 81

\textbf{Obiettivo.} Collegare l’unità di ricampionamento alla gerarchia.

\textbf{Consegna.} Confrontare gli errori standard del bootstrap per righe e per stazioni. Mostrare l’equivalenza della stima osservata nel disegno bilanciato e discutere le ipotesi di indipendenza.

\textbf{Soluzione e interpretazione.} La stima osservata è 2,599 g/cm\textsuperscript{3}. I SE sono 0,014 per righe e 0,024 per stazioni. Il secondo schema assume sei medie locali IID; la numerosità ridotta e la selezione dei siti limitano l’inferenza. Con numerosità diseguali occorre ridefinire esplicitamente i pesi.

\textbf{Variazione.} Confrontare percentile e t sulle sei medie, riportando le rispettive ipotesi.

\texttt{laboratori/M05/M05\_N16.R}

\newpage

\subsection*{N17 · Contrasto fra affioramenti}

\textbf{Pagine:} 80

\textbf{Obiettivo.} Conservare gli affioramenti come strati di stazioni.

\textbf{Consegna.} Ricampionare tre medie di stazione entro A e tre entro B. Calcolare il contrasto A meno B e l’intervallo percentile.

\textbf{Soluzione e interpretazione.} Il contrasto è 0,097 g/cm\textsuperscript{3} e il percentile va da 0,042 a 0,150 g/cm\textsuperscript{3}. Ogni replica mantiene l’appartenenza all’affioramento, ma i tre siti per strato offrono una base molto limitata per la valutazione delle code.

\textbf{Variazione.} Invertire il contrasto e verificare che gli estremi cambino segno e ordine.

\texttt{laboratori/M05/M05\_N17.R}

\subsection*{N18 · Precisione Monte Carlo e registrazione}

\textbf{Pagine:} 86–89

\textbf{Obiettivo.} Distinguere B da n e documentare l’esecuzione.

\textbf{Consegna.} Confrontare dodici esecuzioni per B=200,1000,5000. Rappresentare separatamente SE e i due estremi percentili. Registrare versione di R, algoritmo RNG e semi.

\textbf{Soluzione e interpretazione.} L’aumento di B riduce la dispersione numerica fra corse, mentre il campione conserva ottanta individui. La linea di riferimento del SE è il valore condizionato esatto. La variabilità degli estremi non deve essere interpretata come nuova incertezza della popolazione.

\textbf{Variazione.} Ripetere l’analisi con 24 corse per B, conservando in tabella anche il seme di ciascuna corsa.

\texttt{laboratori/M05/M05\_N18.R}

\newpage

\subsection*{N19 · Reinserimento e bootstrap enumerato}

\textbf{Pagine:} 60–61

\textbf{Obiettivo.} Verificare esattamente il procedimento su un campione piccolo.

\textbf{Consegna.} Enumerare tutte le 256 sequenze ordinate estratte con reinserimento da 4,5,7,8. Calcolare i momenti delle medie e confrontarli con quelli ottenuti permutando tutti e quattro i valori.

\textbf{Soluzione e interpretazione.} La media condizionata è 6 e la varianza condizionata è 0,625, uguale a 3*var(piccolo)/16. Le permutazioni complete conservano la media 6, per cui la loro varianza è zero. I due procedimenti generano distribuzioni differenti.

\textbf{Variazione.} Ripetere il calcolo esatto con il campione 1,1,2,10 e discutere l’effetto del valore elevato.

\texttt{laboratori/M05/M05\_N19.R}

\subsection*{N20 · Mediana e confronto fra statistiche}

\textbf{Pagine:} 68

\textbf{Obiettivo.} Descrivere la distribuzione bootstrap di una statistica non lineare.

\textbf{Consegna.} Calcolare mediana osservata, distribuzione bootstrap e IC percentile per Rounded. Confrontare la forma con quella della distribuzione delle medie.

\textbf{Soluzione e interpretazione.} La mediana osservata è 6,270 µm e l’IC va da 5,957 a 6,452 µm. La distribuzione bootstrap può presentare gradini o concentrazioni perché la mediana dipende dall’ordinamento e può assumere valori ripetuti. Piccoli campioni e distribuzioni irregolari possono compromettere la copertura.

\textbf{Variazione.} Eseguire il bootstrap della mediana sul campione 4,5,7,8 e confrontare la discrezione del risultato.

\texttt{laboratori/M05/M05\_N20.R}

\end{document}